\documentclass[11pt]{article}

\usepackage[T1]{fontenc}
\usepackage[margin=1in]{geometry}
\usepackage{amsmath}
\usepackage{amssymb}
\usepackage{array}
\usepackage{booktabs}
\usepackage{enumitem}
\usepackage{longtable}
\usepackage{xcolor}
\usepackage[hidelinks]{hyperref}

\title{VBOGS Project Report}
\author{VBOGS}
\date{May 18, 2026}

\newcommand{\R}{\mathbb{R}}
\newcommand{\Umax}{U_{\max}}

\setlist[itemize]{leftmargin=*, itemsep=0.2em, topsep=0.3em}
\setlist[enumerate]{leftmargin=*, itemsep=0.2em, topsep=0.3em}
\renewcommand{\arraystretch}{1.18}

\begin{document}
\maketitle

\begin{abstract}
VBOGS is a mid-implementation research system for uncertainty-aware
autonomous-vehicle mapping. It combines Octree-AnyGS, a scalable
level-of-detail Gaussian Splatting scene representation, with Variational
Bayes Gaussian Splatting (VBGS), a Bayesian Gaussian mixture model that
produces posterior uncertainty. The result is a scene model with one
uncertainty value per Octree-AnyGS anchor. Those anchor uncertainties can be
rendered from candidate camera poses to choose a next-best view (NBV): the
view expected to observe the most uncertain visible scene content.

The repository now has owned entry points for the core pipeline from data
preparation through NBV scoring. The remaining work is primarily validation:
checking per-anchor posteriors, running the full pipeline on a scene whose
uncertainty behavior is understood, confirming rendered uncertainty overlays,
and documenting failure modes before treating the outputs as scientific
results.
\end{abstract}

\tableofcontents
\newpage

\section{Project Goal}

The project asks a practical mapping question: after a vehicle has built a
Gaussian Splatting scene from camera data, where should it look next to reduce
uncertainty? Standard Gaussian Splatting methods can render high-quality
images, but they do not directly expose calibrated uncertainty over local
scene structure. VBOGS adds that uncertainty layer.

The key output is a next-best-view camera pose. A candidate pose is good when
it sees scene content whose local posterior uncertainty is high. This is
different from simply choosing a view with many visible pixels: VBOGS scores
uncertainty per unit of visible surface, so the selected view is biased toward
poorly modeled surfaces rather than large already-confident areas.

\section{System Summary}

\subsection{What VBOGS Combines}

\begin{longtable}{@{}p{0.24\linewidth}p{0.68\linewidth}@{}}
\toprule
\textbf{Component} & \textbf{Role in the project} \\
\midrule
\endhead
Octree-AnyGS &
The primary scene model. It stores a multi-level set of Gaussian anchors and
renders the scene with level-of-detail selection. It runs in the PyTorch
environment. \\
VBGS &
The uncertainty model. It fits Bayesian Gaussian mixtures to normalized
position-color samples and returns posterior parameters. It runs in the JAX
environment. \\
KITTI-360 stereo data &
The initial autonomous-driving data source. Perspective stereo frames provide
RGB observations, ground-truth poses, calibration, and stereo pairs for point
cloud extraction. \\
Repo-owned wrappers &
Scripts under \path{scripts/} and helpers under \path{vbogs/} connect the two
upstream systems without modifying the upstream submodules. \\
\bottomrule
\end{longtable}

\subsection{High-Level Data Flow}

\begin{enumerate}
  \item Prepare KITTI-360 frames into the COLMAP-style layout expected by
        Octree-AnyGS.
  \item Train Octree-AnyGS on posed RGB images to obtain anchors, levels,
        opacity, geometry, and appearance parameters.
  \item Estimate stereo depth, unproject pixels, and save a world-frame
        point cloud with RGB values.
  \item Bucket stereo points into Octree-AnyGS anchors using the same grid
        discretization that created the anchors.
  \item Normalize the point data and fit one VBGS posterior per sufficiently
        observed anchor.
  \item Reduce each posterior to one scalar uncertainty value.
  \item Render scalar uncertainty from candidate camera poses and choose the
        highest-scoring NBV candidate.
\end{enumerate}

\section{Core Concepts}

\subsection{Anchors and Voxels}

In this repository, an Octree-AnyGS anchor is treated as a voxel at a specific
level of detail. There is no separate ``leaf voxel'' object. Octree-AnyGS stores
anchors in a flat tensor, and each anchor has a level \(\ell\). Its cell size is

\[
  cell\_size(\ell)
  =
  \frac{\texttt{voxel\_size}}{\texttt{fork}^{\ell}}.
\]

When this report says ``voxel,'' it means ``anchor cell at that level.'' This
distinction matters because anchors from multiple levels coexist. A far camera
may render a coarse anchor while a nearby camera may render finer anchors in the
same physical region.

\subsection{Per-Anchor Uncertainty}

For every anchor \(a\), VBOGS wants a scalar uncertainty \(U_a\). This value is
derived from the VBGS posterior fitted to the stereo points that fall inside the
anchor cell. Observed anchors receive a posterior-entropy score. Anchors with
too few points receive \(\Umax\), a sentinel high-uncertainty value.

The initial uncertainty definition is the expected mixture-weighted entropy of
the per-component posteriors:

\[
  U_a
  =
  \sum_{k=1}^{K_a}
  \bar{\pi}_{ak}
  \left(
    H_{\mathrm{NW}}(q^s_{ak})
    +
    H_{\Delta}(q^c_{ak})
  \right),
\]

where \(\bar{\pi}_{ak}\) is the expected Dirichlet mixture weight,
\(H_{\mathrm{NW}}\) is the spatial Normal-Wishart posterior entropy, and
\(H_{\Delta}\) is the color/delta posterior entropy.

\subsection{Coordinate Frames}

The same stereo point data is used in two coordinate systems:

\begin{itemize}
  \item \texttt{points\_world}: world-frame positions and RGB. These are used
        for anchor bucketing because Octree-AnyGS anchors live in world space.
  \item \texttt{points\_norm}: globally normalized position and RGB. These are
        used for VBGS fitting because the VBGS priors expect normalized inputs.
\end{itemize}

Mixing these frames is one of the main ways to produce plausible-looking but
incorrect results. In particular, bucketing must not use normalized
coordinates.

\section{Architecture and Environments}

VBOGS deliberately keeps the PyTorch and JAX sides separate. Octree-AnyGS and
VBGS have incompatible CUDA dependency stacks, so the project uses two
environments:

\begin{longtable}{@{}p{0.24\linewidth}p{0.68\linewidth}@{}}
\toprule
\textbf{Environment} & \textbf{Responsibilities} \\
\midrule
\endhead
\texttt{vbogs-torch} &
Octree-AnyGS training and rendering, stereo point extraction, point-to-anchor
bucketing, uncertainty rendering, NBV scoring, and most operator-facing
pipeline orchestration. \\
\texttt{vbogs-jax} &
VBGS posterior fitting and posterior-to-uncertainty computations. \\
\bottomrule
\end{longtable}

The framework boundary is the filesystem. PyTorch-side scripts write
\texttt{.npz}, \texttt{.npy}, \texttt{.json}, and visualization artifacts. JAX-side
scripts read those files, produce posterior files, and write uncertainty
arrays back to disk. This avoids in-process tensor sharing and keeps each
framework inside the environment it supports.

The upstream directories \path{Octree-AnyGS/} and \path{vbgs/} are treated as
read-only dependencies. VBOGS functionality lives in repo-owned scripts and
library helpers.

\section{Pipeline Stages}

\subsection{Stage 1: Prepare Data and Train Octree-AnyGS}

The source data is KITTI-360 perspective stereo. The preparation step adapts
KITTI-360 into the COLMAP-style layout expected by Octree-AnyGS:

\begin{itemize}
  \item camera intrinsics and extrinsics;
  \item posed left-camera RGB images;
  \item sparse reconstruction sidecars required by the training code;
  \item metadata that records the source drive and frame selection.
\end{itemize}

Octree-AnyGS is then trained normally from posed RGB frames. Stereo is not used
as a training loss for this stage; it is used later to generate the point
observations that feed the uncertainty head.

The important output is an Octree-AnyGS run directory containing the anchor
point cloud and model parameters. Those anchors define the grid that later
receives uncertainty values.

\subsection{Stage 2: Build a Stereo Point Cloud}

For each stereo pair, VBOGS estimates disparity, converts disparity to depth,
unprojects valid depth pixels into the left-camera frame, transforms those
points into the world frame, attaches RGB from the left image, and concatenates
all frames.

The baseline matcher is OpenCV \texttt{StereoSGBM}. The interface is
provider-agnostic, so future matchers such as RAFT-Stereo can be introduced
without changing downstream file formats.

The output contract is \path{points_world.npz}, with keys:

\begin{longtable}{@{}p{0.24\linewidth}p{0.68\linewidth}@{}}
\toprule
\textbf{Key} & \textbf{Meaning} \\
\midrule
\endhead
\texttt{xyz} & World-frame point positions. \\
\texttt{rgb} & RGB values sampled from the left image. \\
\texttt{frame\_id} & The source KITTI-360 frame id for each point. \\
\bottomrule
\end{longtable}

\subsection{Stage 3: Bucket Points into Anchors}

The bucketing stage assigns world-frame points to Octree-AnyGS anchors. It must
match the discretization used by Octree-AnyGS when anchors were created. For an
anchor level \(\ell\),

\[
  grid\_coord(p,\ell)
  =
  \mathrm{round}
  \left(
    \frac{p - \texttt{init\_pos}}{cell\_size(\ell)}
  \right).
\]

The implementation builds a lookup from \((level, grid\_coord)\) to
\texttt{anchor\_id}, then assigns each point to every matching level. Assigning
only to the finest matching level would make coarse anchors look artificially
uncertain, even in areas with good observations.

The stage also normalizes the point data for VBGS and writes:

\begin{itemize}
  \item \path{points_norm.npz}: normalized \(\texttt{xyz}+\texttt{rgb}\) rows;
  \item \path{pts_by_anchor.npz}: packed per-anchor point indices;
  \item \path{norm_params.json}: global normalization parameters;
  \item \path{bucket_metadata.json}: counts and diagnostics.
\end{itemize}

\subsection{Stage 4: Fit VBGS Posteriors Per Anchor}

For each anchor with at least \texttt{MIN\_POINTS\_PER\_ANCHOR} assigned
points, VBOGS fits a VBGS mixture over normalized six-dimensional samples:

\[
  x_n =
  \begin{bmatrix}
    p_n \\
    c_n
  \end{bmatrix}
  \in \R^6,
  \qquad
  p_n \in \R^3,
  \qquad
  c_n \in \R^3.
\]

The component count starts at \texttt{K\_INIT}. Larger models are accepted only
when their mean per-point evidence lower bound (ELBO) improves by at least
\texttt{ELBO\_IMPROVEMENT\_TOL}. The initial project values are:

\begin{longtable}{@{}p{0.32\linewidth}p{0.60\linewidth}@{}}
\toprule
\textbf{Parameter} & \textbf{Initial value} \\
\midrule
\endhead
\texttt{K\_INIT} & \texttt{pc.n\_offsets = 10}. \\
\texttt{K\_MAX} & \texttt{4 * K\_INIT = 40}. \\
\texttt{K\_GROWTH\_FACTOR} & \texttt{2}, so \(K \mapsto 2K\). \\
\texttt{MIN\_POINTS\_PER\_ANCHOR} & \texttt{20}. \\
\texttt{ELBO\_IMPROVEMENT\_TOL} & \texttt{0.01} nats per point. \\
\bottomrule
\end{longtable}

The posterior output contains fields such as Dirichlet \(\alpha\), final
component count, observed-anchor mask, Normal-Wishart spatial parameters, and
delta/color posterior parameters.

\subsection{Stage 5: Compute Scalar Uncertainty}

The posterior-to-scalar stage computes one value \(U_a\) for every anchor:

\begin{itemize}
  \item observed anchors use the mixture-weighted posterior entropy definition;
  \item unobserved anchors receive \(\Umax\);
  \item metadata and histogram outputs are written for inspection.
\end{itemize}

The main output is \path{U.npy}, a one-dimensional float array with shape
\([N_{\mathrm{anchors}}]\). The \(i\)-th entry corresponds to the \(i\)-th
Octree-AnyGS anchor.

\subsection{Stage 6: Render Uncertainty and Score NBV Candidates}

Octree-AnyGS normally renders learned color through its Gaussian generation
path. VBOGS reuses that same geometry and visibility path, but substitutes the
per-anchor uncertainty value as the rendered scalar color. This is implemented
as \texttt{render\_scalar}.

For each candidate camera pose, scalar rendering returns:

\begin{itemize}
  \item an uncertainty image \(I_U\), whose pixels contain visible uncertainty;
  \item an alpha image \(I_\alpha\), whose pixels indicate visible scene
        support.
\end{itemize}

The NBV score is

\[
  score(cam)
  =
  \frac{\sum_{u,v} I_U(u,v)}
       {\sum_{u,v} I_\alpha(u,v) + \epsilon}.
\]

The alpha-normalized score means ``prefer views of uncertain visible content,''
not simply ``prefer views with the largest amount of visible geometry.'' The
current candidate-source interface supports simple local/test/lattice sources
and is designed to accept planner-produced reachable poses later.

\subsection{Stage 7: Bundle, Inspect, and Validate}

The pipeline can collect curated outputs under \path{outputs/v1_0/<drive>/} and
package them into a zip file. Bundles include prepared-data metadata,
diagnostic point clouds, uncertainty maps, rendered views, NBV diagnostics, and
a run manifest. Bulky checkpoints and full posterior files remain in their
native data locations and are referenced by path.

This stage is not just packaging. Scientific validation requires human review
of the posterior fits, heatmaps, and NBV choices.

\section{Current Progress}

\subsection{Milestone Status}

\begin{longtable}{@{}p{0.16\linewidth}p{0.18\linewidth}p{0.58\linewidth}@{}}
\toprule
\textbf{Milestone} & \textbf{Status} & \textbf{What has been achieved} \\
\midrule
\endhead
M1: Environment setup &
Complete &
Separate Torch and JAX environments are represented by repo-owned helper
scripts and Docker services. Smoke-test commands exist for both upstream
stacks. \\
M2: Train Octree-AnyGS &
Implemented &
KITTI-360 preparation and Octree-AnyGS training wrappers exist. The local
development workflow uses a conservative 16 GB profile; the server budget is
46 GB. Training artifacts are saved under Octree-AnyGS run directories. \\
M3: Stereo to world point cloud &
Implemented &
\path{scripts/stereo_to_pointcloud.py} produces \path{points_world.npz} and
optional PLY diagnostics. The matcher interface currently supports the SGBM
baseline and leaves room for future matchers. \\
M4a: Point to anchor bucketing &
Implemented &
\path{scripts/bucket_points.py} buckets points across all Octree-AnyGS levels,
writes normalized data, and records metadata. A bundled development scene
reported 12,792,935 points, 267,830 anchors across 9 levels, and 104,577
anchors with at least 20 assigned points. \\
M4b: Per-anchor VBGS fit &
Implemented; needs quality pass &
\path{scripts/fit_anchors.py} supports grouped batched fitting with
\texttt{jax.vmap}, deterministic shards, shard merge, and a one-anchor loop
fallback. Smoke artifacts exist, but the full-scene fit still needs
completion and posterior-quality inspection. \\
M5: Posterior to scalar uncertainty &
Implemented &
\path{scripts/compute_uncertainty.py} computes the closed-form entropy
components, writes \path{U.npy}, uncertainty components, metadata, and a
histogram. \\
M6: Scalar rendering and NBV loop &
Implemented; needs full GPU validation &
\path{vbogs/render.py} and \path{scripts/score_nbv.py} implement scalar
rendering, candidate scoring, top-candidate diagnostics, and visualization
support. Syntax and CLI checks exist; a full render validation pass on
completed uncertainty outputs remains. \\
M7: End-to-end validation &
Remaining &
The final validation pass is intentionally human-owned. It requires selecting
a known scene, checking posterior fits, overlaying uncertainty on held-out
views, confirming NBV intuition, and documenting failure modes. \\
\bottomrule
\end{longtable}

\subsection{Implemented Entry Points}

The repository now has direct scripts for the main stages:

\begin{itemize}
  \item \path{scripts/prepare_kitti360_colmap.py}: KITTI-360 to COLMAP-style
        data preparation.
  \item \path{scripts/train_octree_anygs.py}: Octree-AnyGS training wrapper.
  \item \path{scripts/stereo_to_pointcloud.py}: stereo disparity to
        world-frame point cloud.
  \item \path{scripts/bucket_points.py}: point-to-anchor bucketing and data
        normalization.
  \item \path{scripts/fit_anchors.py}: per-anchor VBGS posterior fitting.
  \item \path{scripts/compute_uncertainty.py}: posterior entropy reduction.
  \item \path{scripts/render_uncertainty_views.py}: diagnostic uncertainty
        rendering.
  \item \path{scripts/score_nbv.py}: NBV candidate scoring.
  \item \path{scripts/run_drive_pipeline.py}: orchestrated drive-level
        pipeline.
  \item \path{scripts/bundle_run_outputs.py}: curated artifact bundling.
\end{itemize}

Supporting scripts also exist for global VBGS baselines, VBOGS-vs-VBGS
comparison workflows, experiment packaging, visualization, and output upload.

\subsection{Configuration and Deployment Progress}

The pipeline reads \path{configs/pipeline/default.yaml} by default, with
specialized development and Portainer profiles:

\begin{itemize}
  \item \path{configs/pipeline/dev.yaml}: local Docker Compose development
        stack and smaller hardware settings.
  \item \path{configs/pipeline/portainer.yaml}: server deployment profile with
        external volumes and higher-quality training defaults.
  \item \path{configs/docker/stack.env}: stack-level environment variables,
        image names, autorun controls, upload flags, and transfer-service
        settings.
\end{itemize}

This means the project has moved beyond isolated scripts: it now has an
operator-facing path for repeated drive-level runs.

\section{Artifact Map}

\begin{longtable}{@{}p{0.20\linewidth}p{0.32\linewidth}p{0.40\linewidth}@{}}
\toprule
\textbf{Stage} & \textbf{Typical path} & \textbf{Important contents} \\
\midrule
\endhead
Source data &
\path{data/KITTI-360/} &
Calibration, poses, and rectified stereo images. \\
Prepared data &
\path{/data/COLMAP/<drive>/} &
Images, COLMAP-style camera files, sparse sidecars, metadata. \\
Octree-AnyGS run &
\path{/data/OCTREE-ANYGS/<drive>/<timestamp>/} &
Generated config, input PLY, anchor checkpoint, model outputs. \\
Stereo points &
\path{data/points_world/<drive>/} &
\path{points_world.npz}, optional PLY, metadata. \\
Bucketing and normalization &
\path{data/m4/<drive>/} &
\path{points_norm.npz}, \path{pts_by_anchor.npz},
\path{norm_params.json}, \path{bucket_metadata.json}. \\
Anchor posterior &
\path{data/m4/<drive>/} &
\path{anchor_posterior.npz}, \path{fit_metadata.json}. \\
Uncertainty scalar &
\path{data/m4/<drive>/} &
\path{U.npy}, entropy components, metadata, histogram. \\
Curated outputs &
\path{outputs/v1_0/<drive>/} &
Point clouds, rendered diagnostics, NBV outputs, manifest, packaged zip. \\
\bottomrule
\end{longtable}

\section{Design Decisions Already Made}

\begin{longtable}{@{}p{0.32\linewidth}p{0.60\linewidth}@{}}
\toprule
\textbf{Decision} & \textbf{Current choice} \\
\midrule
\endhead
Stereo data source &
KITTI-360 perspective stereo. The original plan selected
\path{2013_05_28_drive_0008_sync}; current configs and examples may target
other KITTI-360 drives for experiments. \\
Baseline stereo matcher &
OpenCV \texttt{StereoSGBM}, with a provider interface for future matchers. \\
Training budget &
46 GB maximum for server training and inference steps, with conservative local
development settings. \\
NBV candidate source &
A planner-compatible interface. Initial runs can use local/test/lattice
candidates; later runs can pass planner-produced reachable poses. \\
Uncertainty scalar &
Mixture-weighted per-component posterior entropy. The initial implementation
does not use total mixture entropy. \\
Score semantics &
Alpha-normalized expected uncertainty per visible surface. \\
Submodule policy &
Do not edit \path{Octree-AnyGS/} or \path{vbgs/}; wrap or reuse them from
repo-owned code. \\
\bottomrule
\end{longtable}

\section{Validation and Testing Status}

The repository contains focused tests for core helpers and pipeline scripts,
including bucketing, uncertainty computation, rendering helpers, pipeline
orchestration, bundling, upload support, baseline scripts, and comparison
workflows. These tests are important for regression control, but they do not
replace visual and statistical validation of uncertainty quality.

The remaining validation work is:

\begin{enumerate}
  \item Inspect representative anchor posteriors after M4b. A tight,
        well-observed anchor should have low entropy; a sparse, noisy, or
        ambiguous anchor should have high entropy.
  \item Confirm that scalar uncertainty histograms have plausible spread and
        are not collapsed or uniform.
  \item Render uncertainty heatmaps on held-out views and compare them to
        scene intuition.
  \item Score NBV candidates and check whether the winning pose points toward
        known uncertain regions.
  \item Record failure modes before treating the pipeline as validated.
\end{enumerate}

\section{Known Limitations and Risks}

\subsection{Empty-Region Blindness}

\texttt{render\_scalar} only renders through existing Octree-AnyGS anchors.
Truly unseen empty space has no Gaussian and therefore contributes nothing to
the NBV score. The current method chooses views of uncertain known surfaces; it
does not yet solve full exploration of never-observed volume.

\subsection{ELBO-Based Model Selection Is Pragmatic}

The current component-growth rule compares mean per-point ELBO across
component counts. This is useful and simple, but the KL term changes with
model size, so ELBO across \(K\) is not a perfectly unbiased model-selection
criterion. Held-out log likelihood or BIC would be cleaner future options.

\subsection{Stereo Quality Controls Uncertainty Quality}

Poor disparity estimates create noisy or missing points. Textureless regions,
reflective surfaces, distant geometry, thin structures, and calibration errors
can all distort anchor posteriors. Improving stereo confidence filtering or
using a stronger stereo backend may directly improve uncertainty estimates.

\subsection{Large Anchor Counts Are Expensive}

Octree-AnyGS scenes can contain \(10^5\) to \(10^6\) anchors. The JAX batched
path is implemented to control this cost, but full-scene fitting still needs
careful memory and runtime monitoring.

\subsection{Score Semantics May Need Planner Tuning}

The current score selects uncertain content per visible surface. A planner that
wants maximum aggregate information gain may prefer an unnormalized or
travel-cost-adjusted objective. The current interface leaves room for those
extensions.

\section{Future Plan}

\subsection{Immediate Next Steps}

\begin{enumerate}
  \item Complete a full-scene M4b posterior fit on the selected KITTI-360
        drive using the batched JAX path.
  \item Run the human posterior validation pass before relying on M5 outputs.
  \item Compute \path{U.npy} and inspect histograms, metadata, and a sample of
        high/low-uncertainty anchors.
  \item Run full Torch/GPU scalar render validation with completed uncertainty
        values.
  \item Execute M7 on a scene whose uncertain regions are already understood.
  \item Save qualitative overlays, NBV diagnostics, and written failure-mode
        notes with the run bundle.
\end{enumerate}

\subsection{Algorithmic Improvements}

\begin{itemize}
  \item Replace or supplement ELBO-based \(K\) selection with held-out log
        likelihood, BIC, or a calibrated validation split.
  \item Add an optional empty-space uncertainty term, such as an occupancy
        prior or unknown-ray penalty, if the project needs exploration of
        never-seen volume.
  \item Integrate a stronger stereo backend while preserving the existing
        \path{points_world.npz} contract.
  \item Calibrate uncertainty values against held-out geometry or image-space
        reconstruction errors.
  \item Add online updates that reuse VBGS sufficient statistics as new frames
        arrive.
  \item Add motion cost, reachability, field-of-view constraints, and planner
        feasibility checks to NBV scoring.
\end{itemize}

\subsection{Engineering Improvements}

\begin{itemize}
  \item Standardize experiment suites across several KITTI-360 drives.
  \item Preserve exact generated configs, Docker image digests, run manifests,
        and data hashes for reproducibility.
  \item Expand smoke tests for GPU rendering and batched JAX fitting where CI
        hardware permits.
  \item Build compact review artifacts for every full run: uncertainty
        histograms, anchor-fit panels, held-out overlays, top-K NBV views, and
        a short validation note.
  \item Keep upstream submodules clean and continue placing VBOGS behavior in
        \path{scripts/} and \path{vbogs/}.
\end{itemize}

\section{How New Readers Should Approach the Repository}

New readers should start with the conceptual documents before running code:

\begin{enumerate}
  \item \path{docs/manuscript/Algorithm.tex}: authoritative algorithm
        specification.
  \item \path{PLAN.md}: milestone checklist and current implementation status.
  \item \path{docs/documentation/ALGORITHM_OVERVIEW.md}: shorter narrative
        overview.
  \item \path{docs/reference/artifacts.md}: file contracts and output layout.
  \item \path{docs/reference/commands.md}: common commands for running stages.
\end{enumerate}

The most important mental model is that VBOGS is not a replacement for
Octree-AnyGS. Octree-AnyGS remains the renderer and scene representation.
VBOGS adds a Bayesian uncertainty head aligned to Octree-AnyGS anchors, then
uses the renderer to project that uncertainty into candidate views.

\section{Conclusion}

VBOGS has reached the point where the core end-to-end machinery exists in the
repository: data preparation, Octree-AnyGS training wrappers, stereo point
cloud generation, multi-level anchor bucketing, batched per-anchor VBGS
fitting, scalar uncertainty computation, uncertainty rendering, NBV scoring,
and curated artifact bundling. The remaining work is less about inventing the
pipeline and more about proving that its outputs are trustworthy.

The next milestone is therefore validation. Once representative posterior fits,
uncertainty overlays, and NBV selections agree with scene intuition, the
project can move from implementation-complete to experimentally credible. From
there, the most important extensions are stronger model selection, better
stereo, planner-aware candidate scoring, and a principled treatment of
unobserved empty space.

\end{document}
